\documentclass[10pt,a4paper]{amsart}
\usepackage[margin=19mm]{geometry}
\usepackage{amsmath,amssymb,mathtools}
\usepackage[scr=rsfs,cal=euler]{mathalfa}
\usepackage{xfrac}
\usepackage[unicode,hidelinks,bookmarks=false]{hyperref}
\newcommand{\ie}{i.e.\ }
\newcommand{\cf}{cf.\ }
\newcommand{\resp}{respectively\ }
\newcommand{\Subsection}{\subsection}
\providecommand{\nobreakdash}{\nobreak\hbox{-}}
\setlength{\emergencystretch}{2em}
\multlinegap0pt
\usepackage{amssymb}
\usepackage[shortlabels]{enumitem}
\usepackage{xcolor}
\usepackage{tikz}
\newtheorem{lemma}{Lemma}
\newcommand{\arw}{\longrightarrow}
\newcommand{\shO}{\mathcal{O}}
\title{Lagrangian Intersections, Symplectic Reduction and Kirwan Surjectivity}
\author{Naichung Conan Leung, Ying Xie, Yu Tung Yau}
\date{Source: arXiv:2602.11718v1 (CC BY 4.0)}
\begin{document}
\maketitle

\noindent\emph{Source and licence.} Lagrangian Intersections, Symplectic Reduction and Kirwan Surjectivity. Version arXiv:2602.11718v1 (CC BY 4.0), distributed by arXiv under the Creative Commons Attribution 4.0 International licence, http://creativecommons.org/licenses/by/4.0/, as recorded in the arXiv licence metadata for this exact version and shown on the abstract page https://arxiv.org/abs/2602.11718v1. Full licence notice: \texttt{licenses/arxiv-2602.11718-CC-BY-4.0.txt}.\par\smallskip

\noindent\textit{Editorial scope.} Counted unit: the article's lemma canlag (source0.tex lines 536--542). The article's complete original proof of it is delivered with it (source0.tex lines 543--557, one \texttt{proof} environment of 15 lines). No mathematics is written, repaired or re-derived: the statement and the proof are the article's own lines, in the article's own notation, and no retained line is substituted or re-flowed.  Retained uncounted as the source-local context the proof needs: lines 409--412; lines 525--527.  Nothing was omitted from any retained span, so the package carries no omission record.  No retained line contains a citation, so the wrapper carries no bibliography and no bibliography entry.  No retained line contains a figure environment, an image-inclusion command or drawing code, so no image is delivered and the package declares no asset dependency.  Editorial formatting only, in addition: an AMS \texttt{amsart} preamble that restates the article's own declarations and macros used by the retained lines, namely the environments \texttt{lemma} on separate counters, together with the standard packages those lines need.  The retained environments print in this document as Lemma 1 in order of appearance, whereas the article prints the same items with the numbering of its own full document.  The complete native source of the article as distributed in the arXiv e-print for this exact version is bundled unchanged as \texttt{sources/194487/source0.tex}.
\begin{lemma}\label{excess}
   There is a short exact sequence of vector bundles on $B$:
$$0\arw N_{B/C_2}\arw N_{C_1/X}|_B\arw E\arw 0.$$
\end{lemma}

\begin{lemma}\label{exc}
    If $C_1$ and $C_2$ are two smooth Lagrangians of $X$ with a clean intersection $B$, then the excess bundle is isomorphic to the cotangent bundle of $B$: $$E\cong  \Omega_{B}.$$
\end{lemma} 

\begin{lemma}\label{lem: canlag}
Let $C_1$ and $C_2$ be two smooth Lagrangians of $X$ with a clean intersection $B$. Then
\begin{align*}
    (K_{C_1}^{\vee}\otimes K_{C_2})|_B\otimes \det(N_{B/C_2})^{\otimes2}\cong \shO_B.
\end{align*}
    where $K_{C_i}$ is the canonical bundle on $C_i$ ($i=1, 2$).
\end{lemma}

\begin{proof}
    By Lemma \ref{excess} and \ref{exc}, there is short exact sequence:
    \begin{align*}
        0\arw N_{B/C_2} \arw N_{C_1/X}|_B\arw \Omega_B\arw 0.
    \end{align*}
    Note that $N_{C_1/X}\cong \Omega_{C_1}$, and then
    \begin{align}\label{eq: can1}
        K_{C_1}|_B\cong K_B\otimes \det(N_{B/C_2}).
    \end{align}
    The determinant of the conormal sequence of $B\subset C_2$ gives
    \begin{align}\label{eq: can2}
        K_{C_2}\cong K_B\otimes \det(N_{B/C_2}^{\vee}).
    \end{align}
    Thus $(K_{C_1}^{\vee}\otimes K_{C_2})|_B\otimes \det(N_{B/C_2})^{\otimes2}\cong \shO_B$ holds by (\ref{eq: can1}) and (\ref{eq: can2}).
\end{proof}

\end{document}
